%%%%%%%%%%%%%%
%
%  第三部分 总结展望
%
%%%%%%%%%%%%%%
\section{总结展望}
\subsection{总结}

\par
"\SoftwareName（\Version）"围绕超声检测 A-Scan 典型流程，完成了仿真采集、波形实时显示、闸门检测、游标读数、缺陷记录与报告追溯等核心能力。软件采用 MVC 架构实现界面与逻辑分离，界面风格为深色工业风，交互方式贴近现场检测习惯，能够支持培训演示、算法验证与软件工程化扩展。
\par
在使用层面，软件强调"可控采集"和"可追溯记录"：通过刷新间隔、冻结/单次采样与触发模式实现可视化观察；通过"开始记录"开关避免误触发写入历史，保证报告结论与数据一致。

\par
在 V1.0 版本基础上，本软件进一步增强了以下能力：
\begin{enumerate}
	\item \textbf{信号处理能力}：集成 FFT 频谱分析功能，支持频率域信号观察与峰值频率识别；提供带通、低通、高通三种数字滤波器，可对波形进行频率选择性处理；支持多次采样信号平均，提高信噪比。
	\item \textbf{数据分析能力}：新增缺陷统计分析模块，提供缺陷总数、深度分布、波幅分布等统计信息；集成 DAC（距离-波幅校正）曲线生成与缺陷超限判定功能；支持信号特征提取与信号对比分析。
	\item \textbf{报告生成能力}：报告系统升级为多标准模板架构，支持通用、NB/T 47013、GB/T 11345、ASME 等检测标准；报告内容支持公司信息、检测人员、检测参数等定制化配置。
	\item \textbf{用户体验优化}：提供完整的快捷键系统（F1$\sim$F6 及 Ctrl 组合键），大幅提升操作效率；实现界面布局自动记忆功能，窗口位置、面板停靠状态关闭时自动保存、启动时自动恢复；新增快捷键帮助对话框，便于用户查阅。
\end{enumerate}




\subsection{展望}

\par
后续可在保持现有架构的基础上，进一步完善与工程应用相关的能力，例如：
\begin{enumerate}
	\item 硬件接入：对接真实超声采集卡/仪器，实现多通道数据采集与设备状态管理；
	\item 多闸门与曲线：增加 Gate B/C、TCG（时间增益补偿）与多参考曲线管理；
	\item 扫查显示：扩展 B-Scan/C-Scan 或 A-Scan 堆叠回放，实现扫查轨迹与缺陷分布可视化；
	\item 数据回放：支持历史数据加载、时序回放、关键帧标注与缺陷复现；
	\item 报告增强：支持 PDF 导出、缺陷统计图表嵌入与附件打包；
	\item 权限与审计：增加用户登录、操作日志与报告签名，满足更严格的质量体系要求。
\end{enumerate}


